\begin{frame}{Measurement framework}
\framesubtitle{Psychometric basis}

\begin{block}{The measurement problem: three requirements}
\begin{enumerate}
    \item \textbf{Target construct}: what is being measured?
    \item \textbf{Measurement model}: how do scores relate to the construct?
    \item \textbf{Validity evidence}: what evidence supports the interpretation?
\end{enumerate}
\end{block}

\vspace{0.3cm}

\begin{itemize}
    \item Evaluation can target \textbf{systems} or \textbf{artefacts}: model capability across tasks vs. quality of an individual text.

    \item Psychometrics distinguishes between \textbf{reflective} and \textbf{formative} measurement models \citep{bollen1991conventional}.

    \item These models define how scores relate to the construct, what the main validity risks are, and what kind of evidence is needed.
\end{itemize}

\end{frame}




\begin{frame}{Measurement framework}
\framesubtitle{Measurement models}

\centering
\resizebox{0.95\textwidth}{!}{%
    \begin{tikzpicture}[
    >=Latex,
    font=\sffamily\small,
    latent/.style={
        circle,
        draw=black!70,
        thick,
        minimum size=1.15cm,
        align=center
    },
    indicator/.style={
        rectangle,
        rounded corners=2pt,
        draw=black!70,
        thick,
        minimum width=1.15cm,
        minimum height=0.65cm,
        align=center
    },
    title/.style={
        font=\sffamily\bfseries\small,
        align=center
    },
    note/.style={
        font=\sffamily\footnotesize,
        align=center,
        text width=5cm
    },
    arrow/.style={
        -{Latex[length=2.3mm]},
        thick
    }
]

% ------------------------------------------------
% Reflective / system-level
% ------------------------------------------------

\node[title] at (0,2.1) {System-level evaluation};
\node[title] at (0,1.65) {Reflective measurement};

\node[latent] (etaS) at (0,0.55) {$\eta_S$};

\node[indicator] (sx1) at (-1.8,-0.85) {$x_1$};
\node[indicator] (sx2) at (0,-1.15) {$x_2$};
\node[indicator] (sx3) at (1.8,-0.85) {$x_3$};

\draw[arrow] (etaS) -- (sx1);
\draw[arrow] (etaS) -- (sx2);
\draw[arrow] (etaS) -- (sx3);

\node[note] at (0,-2.15)
{Latent system ability $\eta_S$ causes observed performance across tasks, prompts, or benchmark items.};

\node[note] at (0,-3.0)
{$x_j = \lambda_j \eta_S + \varepsilon_j$};

% ------------------------------------------------
% Formative / artefact-level
% ------------------------------------------------

\node[title] at (7.1,2.1) {Artefact-level evaluation};
\node[title] at (7.1,1.65) {Formative measurement};

\node[indicator] (ax1) at (5.3,0.55) {$x_1$};
\node[indicator] (ax2) at (7.1,0.85) {$x_2$};
\node[indicator] (ax3) at (8.9,0.55) {$x_3$};

\node[latent] (etaA) at (7.1,-0.85) {$\eta_A$};

\draw[arrow] (ax1) -- (etaA);
\draw[arrow] (ax2) -- (etaA);
\draw[arrow] (ax3) -- (etaA);

\node[note] at (7.1,-2.15)
{Criterion-specific scores jointly constitute the quality $\eta_A$ of an individual artefact.};

\node[note] at (7.1,-3.0)
{$\eta_A = \sum_{j=1}^{J} \gamma_j x_j + \zeta$};

\end{tikzpicture}
}


\end{frame}


\begin{frame}{Measurement framework}
\framesubtitle{Implication for the empirical design}

\begin{itemize}
    \item In formative measurement, the criterion set is part of the construct specification: changing the criteria changes the construct \citep{bollen1991conventional}

    \item Validity evidence is collected by relating LLM-derived measurements to external validation criteria

    \item The empirical question is whether explicit formative specification yields stronger validity evidence and more interpretable diagnostics.
\end{itemize}

\end{frame}